\documentclass[aspectratio=169,xcolor={dvipsnames,table}]{beamer}
\usepackage[no-math,deluxe,haranoaji]{luatexja-preset}
\renewcommand{\kanjifamilydefault}{\gtdefault}
\renewcommand{\emph}[1]{{\upshape\bfseries #1}}
\usetheme{metropolis}
\metroset{block=fill}
\setbeamertemplate{navigation symbols}{}
\setbeamertemplate{blocks}[rounded][shadow=false]
\usecolortheme[rgb={0.7,0.2,0.2}]{structure}
%%%%%%%%%%%%%%%%%%%%%%%%%%
%% Change alert block colors
%%% 1- Block title (background and text)
\setbeamercolor{block title alerted}{fg=mDarkTeal, bg=mLightBrown!45!yellow!45}
\setbeamercolor{block title example}{fg=magenta!10!black, bg=mLightGreen!70}
%%% 2- Block body (background)
\setbeamercolor{block body alerted}{bg=mLightBrown!25}
\setbeamercolor{block body example}{bg=mLightGreen!15}
%%%%%%%%%%%%%%%%%%%%%%%%%%%
\usepackage[absolute,overlay]{textpos}
%\usepackage[grid=true,gridcolor=Maroon,subgridcolor=gray,gridunit=pt,texcoord]{eso-pic} %場所決めのためのgrid表示
%%%%%%%%%%%%%%%%%%%%%%%%%%%
%% さまざまなアイコン
%%%%%%%%%%%%%%%%%%%%%%%%%%%
%\usepackage{fontawesome}
\usepackage{fontawesome5}
\usepackage{figchild}
\usepackage{twemojis}
\usepackage{utfsym}
\usepackage{bclogo}
\usepackage{marvosym}
\usepackage{fontmfizz}
\usepackage{pifont}
\usepackage{phaistos}
\usepackage{worldflags}
\usepackage{jigsaw}
\usepackage{tikzlings}
\usepackage{tikzducks}
\usepackage{scsnowman}
\usepackage{epsdice}
\usepackage{halloweenmath}
\usepackage{svrsymbols}
\usepackage{countriesofeurope}
\usepackage{tipa}
\usepackage{manfnt}
%%%%%%%%%%%%%%%%%%%%%%%%%%%
\usepackage{tikz}
\usetikzlibrary{calc,patterns,decorations.pathmorphing,backgrounds}
\usepackage{tcolorbox}
\usepackage{tikzpeople}
\usepackage{circledsteps}
\usepackage{xcolor}
\usepackage{amsmath}
\usepackage{booktabs}
\usepackage{chronology}
\usepackage{signchart}
%%%%%%%%%%%%%%%%%%%%%%%%%%%
%% 場合分け
%%%%%%%%%%%%%%%%%%%%%%%%%%%
\usepackage{cases}
%%%%%%%%%%%%%%%%%%%%%%%%%%
\usepackage{pdfpages}
%%%%%%%%%%%%%%%%%%%%%%%%%%%
%% 音声リンク表示
\newcommand{\myaudio}[1]{\href{#1}{\faVolumeUp}}
%%%%%%%%%%%%%%%%%%%%%%%%%%
%% \myAnch{<名前>}{<色>}{<テキスト>}
%% 指定のテキストを指定の色の四角枠で囲み, 指定の名前をもつTikZの
%% ノードとして出力する. 図には remember picture 属性を付けている
%% ので外部から参照可能である.
\newcommand*{\myAnch}[3]{%
  \tikz[remember picture,baseline=(#1.base)]
    \node[draw,rectangle,line width=1pt,#2] (#1) {\normalcolor #3};
}
%%%%%%%%%%%%%%%%%%%%%%%%%%
%% \myEmph コマンドの定義
%%%%%%%%%%%%%%%%%%%%%%%%%%
%\newcommand{\myEmph}[3]{%
%    \textbf<#1>{\color<#1>{#2}{#3}}%
%}
\usepackage{xparse} % xparseパッケージの読み込み
\NewDocumentCommand{\myEmph}{O{} m m}{%
    \def\argOne{#1}%
    \ifx\argOne\empty
        \textbf{\color{#2}{#3}}% オプション引数が省略された場合
    \else
        \textbf<#1>{\color<#1>{#2}{#3}}% オプション引数が指定された場合
    \fi
}
%%%%%%%%%%%%%%%%%%%%%%%%%%%
%%%%%%%%%%%%%%%%%%%%%%%%%%%
%% 文末の上昇イントネーション記号 \myRisingPitch
%% 通常のイントネーション \myDownwardPitch
%% https://note.com/dan_oyama/n/n8be58e8797b2
%%%%%%%%%%%%%%%%%%%%%%%%%%%
\newcommand{\myRisingPitch}{
\begin{tikzpicture}[scale=0.3,baseline=0.3]
\draw[->,>=stealth] (0,0) to[bend right=45] (1,1);
\end{tikzpicture}
}
\newcommand{\myDownwardPitch}{
\begin{tikzpicture}[scale=0.3,baseline=0.3]
\draw[->,>=stealth] (0,1) to[bend left=45] (1,0);
\end{tikzpicture}
}
%%%%%%%%%%%%%%%%%%%%%%%%%%%%
%\AtBeginSection[%
%]{%
%  \begin{frame}[plain]\frametitle{授業の流れ}
%     \tableofcontents[currentsection]
%   \end{frame}%
%}

\usepackage{pxrubrica}
\usetikzlibrary{tikzmark}
%%%%%%%%%%%%%%%%%%%%%%%%%%%
\title{English is fun.}
\subtitle{John is as tall as his father.}
\author{}
\institute[]{}
\date[]

%%%%%%%%%%%%%%%%%%%%%%%%%%%%
%% TEXT
%%%%%%%%%%%%%%%%%%%%%%%%%%%%
\begin{document}

\begin{frame}[plain]
  \titlepage
\end{frame}

%%%%%%%%%%%%%%%%%%%%%%%%%%%%%%%
\section*{授業の流れ}
\begin{frame}[plain]
  \frametitle{授業の流れ}
  \tableofcontents
\end{frame}
%%%%%%%%%%%%%%%%%%%%%%%%%%%%%%%
\section{同じくらい～\hspace{80pt}{\tiny 2042 audio overview}\,{\scriptsize \myaudio{./audio/overview/041_as_as_audio_overview.mp4}}
}
\subsection{as ～ as \ldots{}}
%%%%%%%%%%%%%%%%%%%%%%%%%%%%%%%%%%%%%%%%%%%%%
\begin{frame}[plain]{as ～ as \ldots{}}
 %\Large

\begin{enumerate}
 \item \begin{enumerate}
	\item<1-> John is \myEmph[1-]{NavyBlue}{tall}.
	\item<2-> John is \myEmph[2-]{Maroon}{as} \myEmph[2-]{NavyBlue}{tall} \myEmph[2-]{Maroon}{as} his brother.\hfill{\textipa{/@z/}}
       \end{enumerate}
 \item \begin{enumerate}
	\item<3-> Jennifer runs \myEmph[3-]{NavyBlue}{fast}. 
	\item<4-> Jennifer runs \myEmph[4-]{Maroon}{as} \myEmph[4-]{NavyBlue}{fast} \myEmph[4-]{Maroon}{as} Emily.
       \end{enumerate}
\end{enumerate}
%
\hfill{\tiny 0152}\,{\scriptsize \myaudio{./audio/041_as_as_01.mp3}}
\vfill

\begin{block}<5->{Topics for Today}
\begin{itemize}\setbeamertemplate{items}[square]\small
 \item<6-> AとBの程度が同じとき（$\text{A}=\text{B}$のとき）は\\
\mbox{}\hspace{120pt} $\text{\Circled[fill color=white]{A}\,\,\,\,\,\ldots\,\,\,\,\,\myEmph[5-]{Maroon}{as}\,\,}+\left\{\begin{tblr}{l}
	    \text{形容詞}\\
	    \text{副詞}
	 \end{tblr}\right\} + \text{\myEmph[5-]{Maroon}{as}\,\,\,\,\,\Circled[fill color=white]{B}}$
 \item<7-> 述語動詞はbe動詞のこともあれば一般動詞のこともあります
 \item<8-> \kenten{何}と\kenten{何}を\kenten{何}について比べているのか確認しよう\\%
       \hfill\visible<9->{{\scriptsize John($=$ \Circled[fill color=white]{A})とhis brother($=$ \Circled[fill color=white]{B})を「背の高さ」($=$ tall)について比較}}\\
       \hfill\visible<10->{{\scriptsize Jennifer($=$ \Circled[fill color=white]{A})とEmily($=$ \Circled[fill color=white]{B})を「走る速さ」($=$ fast)について比較}}
 \end{itemize}
     \end{block}

\begin{textblock*}{0.4\linewidth}(358pt,130pt)
\visible<11->{\begin{tikzpicture}
\pig[
signpost=\scalebox{0.5}{
\parbox{2.2cm}{\color{black}
\centering 何と何を\\何について\\比べてる？}},
%speech={\tiny
% \begin{tabular}{c}
% 目的を表す\\to不定詞
% \end{tabular}
%},
signcolour= brown!70!gray,
signback=white!80!brown
]
\end{tikzpicture}}
\end{textblock*}

\end{frame}
%%%%%%%%%%%%%%%%%%%%%%%%%%%%%%%%%%%%%%%%%%%%%%
\begin{frame}[plain]{Exercises}
\small
日本語の意味になるよう(~~~~~~)内の語句を並べ替えましょう。なお $[ +1 ]$とある場合は不足している1語を補ってください。$[ -1 ]$とある場合は余計な1語が含まれています%
\hfill{\tiny 0215}\,{\scriptsize \myaudio{./audio/041_as_as_02.mp3}}

\normalsize
 \begin{enumerate}
  \item {\small その国ではバスケットボールは野球と同じくらい人気がある。}\hfill{\scriptsize popular \textipa{/p\'ApjUl\textrhookschwa /} 人気のある}\\
	Basketball is ( popular / as / as ) baseball in the country.\hfill{\scriptsize country \textipa{/k\'\textturnv ntri/} 国}\\
	\hspace{80pt}\visible<2->{\myEmph[2-]{BurntOrange}{as popular as}}
  \item {\small ジョージはジョンと同じくらい有名だ。}\hfill{\scriptsize famous \textipa{/f\'eIm@s/} 有名な}\\
	George ( famous / as / is / as ) John.\\
	\hspace*{55pt}\visible<3->{\myEmph[3-]{BurntOrange}{is as famous as}}
  \item {\small ゴリラは象と同じくらいの速さで走る。}\\
	Gorillas ( run / is / fast / as / as ) elephants. $[-1]$\\
	\hspace{65pt}\visible<4->{\myEmph[4-]{BurntOrange}{run as fast as}　✕ is}
  \item {\small ジェームズはエリックと同じくらいギターをじょうずに弾いた。}\\
	James played the guitar ( as / well ) Eric. $[+1]$\\
	\hspace{125pt}\visible<5->{\myEmph[5-]{BurntOrange}{as well as}}
 \end{enumerate}
\end{frame}
%%%%%%%%%%%%%%%%%%%%%%%
\section{as ～ as \ldots\,\,\,の否定}
\subsection{not $\left\{\begin{tabular}{l}
                    as\\
                    so
		     \end{tabular}\right\}$
		     ～ as \ldots}
%%%%%%%%%%%%%%%%%%%%%%
\begin{frame}[plain,t]{not as ～ as\hspace{10pt} \visible<4->{/\hspace{10pt}not so ～ as}}
 \begin{enumerate}
  \item<1-> Tokyo Tower is \myEmph[1-]{Maroon}{not} \myEmph[1-]{NavyBlue}{as} high \myEmph[1-]{NavyBlue}{as} Tokyo Skytree.\hfill{\scriptsize high \textipa{/h\'aI/} 高い}
  \item<3-> Tokyo Tower is \myEmph[2-]{Maroon}{not} \myEmph[2-]{NavyBlue}{so} high \myEmph[2-]{NavyBlue}{as} Tokyo Skytree.
 \end{enumerate}
%
\vspace*{40pt}

\begin{block}<2->{Topics for Today}
\begin{itemize}\setbeamertemplate{items}[square]\small
 \item<2-> $\text{\Circled[fill color=white]{A}\,\,\,\,\,\ldots\,\,\,\,\,\myEmph[2-]{Maroon}{not}\,\,\,\,\myEmph[2-]{Maroon}{as}\,\,}+\left\{\begin{array}{l}
	    \text{形容詞}\\
	    \text{副詞}
	 \end{array}\right\} + \text{\myEmph[2-]{Maroon}{as}\,\,\,\,\,\Circled[fill color=white]{B}}$%
\,\,\,\,$\Longrightarrow$\,\,\,\,「AはBほど～でない」の意味%
      \hfill{A $<$ B}
 \item<4-> 最初の\textbf{as}が\textbf{so}になることがある\\
\hspace{100pt}%
$\text{\Circled[fill color=white]{A}\,\,\,\,\,\ldots\,\,\,\,\,\myEmph[4-]{Maroon}{not}\,\,\,\,\myEmph[2-]{NavyBlue}{so}\,\,}+\left\{\begin{array}{l}
	    \text{形容詞}\\
	    \text{副詞}
	 \end{array}\right\} + \text{\myEmph[4-]{Maroon}{as}\,\,\,\,\,\Circled[fill color=white]{B}}$%
 \item<5-> \kenten{何}と\kenten{何}を\kenten{何}について比べているのか確認しよう\\%
\hfill\visible<6->{{\scriptsize Tokyo Tower($=$ \Circled[fill color=white]{A})とTokyo Skytree($=$ \Circled[fill color=white]{B})を「建物の高さ」($=$ high)について比較}}
 \end{itemize}
     \end{block}

\vspace{-5pt}

\hfill{\tiny 0128}\,{\scriptsize \myaudio{./audio/041_as_as_03.mp3}}

\end{frame}
%%%%%%%%%%%%%%%%%%%%%%%
\begin{frame}[plain]{Exercises}\small
日本語の意味になるよう(~~~~~~)内の語句を並べ替えましょう%
\hfill{\tiny 0146}\,{\scriptsize \myaudio{./audio/041_as_as_04.mp3}}

\normalsize
 \begin{enumerate}
  \item {\small わたしはピーターほど忙しくありません。}%
\hfill{\scriptsize buzy \textipa{/b\'Izi/} 忙しい }\\
	I am ( busy / as / as / not )  Peter. \\
	\hspace{45pt}\visible<2->{\myEmph[2-]{BurntOrange}{not as busy as}}
  \item {\small その映画はその本ほどおもしろくない。}\hfill{\scriptsize interesting \textipa{/\'Intr@stiN/} おもしろい}\\
	The movie is ( as / so / not / interesting ) the book.\\
	\hspace{75pt}\visible<3->{\myEmph[3-]{BurntOrange}{not so interesting as}}
  \item {\small 彼女はジェームズほどピアノの演奏がうまくない。}\\
	She does not play the piano ( as / as / well ) James.\\
	\hspace{160pt}\visible<3->{\myEmph[3-]{BurntOrange}{as well as}}
 \end{enumerate}
\end{frame}
%%%%%%%%%%%%%%%%%%%%%%%%%
\section{読んでみよう}
%%%%%%%%%%%%%%%%%%%%%%
\begin{frame}[plain,t]

\small
\begin{tcolorbox}\parindent=15pt
\noindent{}Today, our class has a lunch contest.%
\hfill{\tiny 0033}\,{\scriptsize \myaudio{./audio/041_as_as_05.mp3}}

    First, we look at the sandwiches. Tom's sandwich is big, and it's as big\tikzmark{clue1} as Mike's. Both are very big!

    Next, we look at the salads and fruit. Sally's \tikzmark{clue2}salad has many colors. It's as colorful as Anna's fruit bowl. Anna's fruit is also very colorful. Ken has soup. It's very hot. His soup is as \tikzmark{clue3}hot as a cup of tea.

    Then, we start eating. Tom eats his sandwich as \tikzmark{clue4}fast as Mike.  Sally eats her salad slowly. She eats as slowly as Anna.

    After we all finish, we decide on the winner. We all agree. Every lunch is a winner!% Each lunch is as good as the others.
\end{tcolorbox}

\visible<2->{{\small つぎの英文が本文の内容に一致すればT、
一致しなければFと答えましょう}}

\begin{enumerate}\small\setlength{\itemsep}{01pt}
 \item<2-> Tom's sandwich is not so big as Mike's.\tikzmark{q1}\hfill\visible<4->{F}
 \item<2-> Anna's salad is as colorful as Sally's fruit bowl.\tikzmark{q2}\hfill\visible<6->{F}
 \item<2-> Ken's soup is as hot as the tea.\tikzmark{q3}\hfill\visible<8->{T}
 \item<2-> Tom eats as quickly as Mike.\tikzmark{q4}\hfill\visible<10->{T}
\end{enumerate}

\begin{tikzpicture}[remember picture,overlay]
 \visible<3->{\draw[<-,opacity=.4,line width=2pt] ([xshift=-6pt,yshift=-2pt]pic cs:clue1) to[bend left] ([xshift=2pt, yshift=2pt] pic cs:q1);}
 \visible<5->{\draw[<-,opacity=.4,line width=2pt] ([yshift=-2pt]pic cs:clue2) to[bend left] ([xshift=2pt, yshift=2pt] pic cs:q2);}
 \visible<7->{\draw[<-,line width=2pt, opacity=.4] ([xshift=8pt,yshift=-2pt]pic cs:clue3) to[bend left] ([xshift=5pt, yshift=2pt] pic cs:q3);}
 \visible<9->{\draw[<-,opacity=.4,line width=2pt] ([xshift=5pt,yshift=-2pt]pic cs:clue4) to[bend left] ([xshift=2pt, yshift=2pt] pic cs:q4);}
\end{tikzpicture}

\normalsize
\begin{textblock*}{0.4\linewidth}(275pt,180pt)
\visible<2->{\begin{tikzpicture}
\duck[signpost=\scalebox{0.3}{
\parbox{2.5cm}{\color{black}
{\Large T or F?}}},
signcolour=brown!70!gray,
signback=white!80!brown,
graduate=gray!20!black,
tassel=red!70!black,
speech={\tiny 根拠をはっきり}
]
\end{tikzpicture}}
\end{textblock*}
\end{frame}
%%%%%%%%%%%%%%%%%%%%%%%%%
\begin{frame}[plain]{大意}
 
今日は、私たちのクラスでランチコンテストがあります。

最初に、サンドイッチを見てみましょう。トムのサンドイッチは大きいです。そして、それはマイクのサンドイッチと同じくらい大きいです。どちらもとても大きいです！

次に、サラダとフルーツを見ます。サリーのサラダはたくさんの色があります。それはアンナのフルーツボウルと同じくらいカラフルです。アンナのフルーツもとてもカラフルです。ケンはスープを持っています。それはとても熱いです。彼のスープは一杯のお茶と同じくらい熱いです。

それから、私たちは食べ始めます。トムはマイクと同じくらい速くサンドイッチを食べます。サリーはサラダをゆっくり食べます。彼女はアンナと同じくらいゆっくり食べます。

みんなが食べ終わった後、私たちは勝者を決めます。みんな意見が一致します。どのランチも勝者です！
\end{frame}
%%%%%%%%%%%%%%%%%%%%%%%%%%
\section{as ～ as \ldots\,\,\,のまとめ}
%%%%%%%%%%%%%%%%%%%%%%%%%
\begin{frame}[plain]{まとめ1}
 \begin{block}<1->{as ～ as \ldots}
\begin{itemize}\setbeamertemplate{items}[square]\small
 \item<1-> AとBの程度が同じとき（$\text{A}=\text{B}$のとき）は\\
\mbox{}\hspace{120pt} $\text{\Circled[fill color=white]{A}\,\,\,\,\,\ldots\,\,\,\,\,\textbf{as}\,\,}+\left\{\begin{tblr}{l}
	    \text{形容詞}\\
	    \text{副詞}
	 \end{tblr}\right\} + \text{\textbf{as}\,\,\,\,\,\Circled[fill color=white]{B}}$
 \item<2-> 述語動詞はbe動詞のこともあれば一般動詞のこともあります
 \item<3-> \kenten{何}と\kenten{何}を\kenten{何}について比べているのか確認しよう\\%
 \end{itemize}
     \end{block}

\visible<4->{\begin{enumerate}
 \item This cake is \textbf{as} sweet \textbf{as} honey.\hfill{\scriptsize honey \textipa{/h\'2ni/} はちみつ}
 \item Eating a salad is \textbf{as} good \textbf{as} eating an apple.\hfill{\scriptsize salad \textipa{/s\'\ae l@d/} サラダ}
 \item She sings \textbf{as} well \textbf{as} her mother.
 \item We drink tea \textbf{as} often \textbf{as} coffee.\hfill{\scriptsize often \textipa{/\'O:f(t)n/} しばしば}
\end{enumerate}}

\hfill{\tiny 0209}\,{\scriptsize \myaudio{./audio/041_as_as_06.mp3}}

\end{frame}
%%%%%%%%
%%%%%%%%%%%%%%%%%%%%%%%%%
\begin{frame}[plain]{まとめ2}
\begin{block}<1->{as ～ as \ldots\,\,の否定}
\begin{itemize}\setbeamertemplate{items}[square]\small
 \item<1-> $\text{\Circled[fill color=white]{A}\,\,\,\,\,\ldots\,\,\,\,\,\textbf{not}\,\,\,\,\textbf{as}\,\,}+\left\{\begin{array}{l}
	    \text{形容詞}\\
	    \text{副詞}
	 \end{array}\right\} + \text{\textbf{as}\,\,\,\,\,\Circled[fill color=white]{B}}$%
\,\,\,\,$\Longrightarrow$\,\,\,\,「AはBほど～でない」の意味%
      \hfill{A $<$ B}
 \item<2-> 最初の\textbf{as}が\textbf{so}になることがある\\
\hspace{100pt}%
$\text{\Circled[fill color=white]{A}\,\,\,\,\,\ldots\,\,\,\,\,\textbf{not}\,\,\,\,\myEmph[2-]{Maroon}{so}\,\,}+\left\{\begin{array}{l}
	    \text{形容詞}\\
	    \text{副詞}
	 \end{array}\right\} + \text{\textbf{as}\,\,\,\,\,\Circled[fill color=white]{B}}$%
 \end{itemize}
     \end{block}

\visible<3->{\begin{enumerate}
 \item My house is \textbf{not as} old \textbf{as} his.\hfill{\scriptsize his $=$ his house}
 \item This car is \textbf{not so} popular \textbf{as} that one.\hfill{\scriptsize that one $=$ that car}
 \item She does \textbf{not} eat \textbf{as} quickly \textbf{as} Bob.
 \item He does \textbf{not} enjoy pizza  \textbf{so} much \textbf{as} pasta.\hfill{\tiny 0215}\,{\scriptsize \myaudio{./audio/041_as_as_07.mp3}}

\end{enumerate}}


\hfill{\tiny audio\_overview 1305}\,{\scriptsize \myaudio{./audio/overview/041_as_as_audio_overview.mp4}}

\end{frame}
%%%%%%%%
\end{document}
